% status: ZERO
% Chapter 21 of the second edition. Plan: docs/STRUCTURE.md. Rules: AUTHORING.md.
% Measurements: research/measurements/2026-09-24-m1-part3.md (lora, ollama_switch).
\chapter{One Server, Many Models: LoRA, Adapters and Routing}\label{ch:multi-model-serving}

\begin{ChapterIntent}
Fine-tuning produces thousands of model variants that differ from a base
model by small low-rank adapters. Serving each on its own GPU wastes almost all
of it. This chapter shows how one server batches requests that use different
adapters --- gathering the right adapter weights per request inside the kernel
--- how adapters are paged in and out of accelerator memory, and how routing
and model multiplexing let a fleet serve many models with little duplicated
capacity.
\end{ChapterIntent}

\OpeningSection{A thousand fine-tunes, one GPU}

A rank-16 LoRA adapter on every projection of Llama~3.2~3B holds 24.3 million
parameters: 48.6~MB in 16-bit floats, less than 1\% of the model's 3.2 billion
parameters (derived). A thousand such fine-tunes fit in 49~GB of host memory;
a thousand fully fine-tuned copies would need 6.4~TB in BF16. On paper, one
server can hold them all.

To see what serving them costs, we loaded three synthetic adapters into
llama.cpp's server on the M1 --- two of rank 16, one of rank 64 --- and chose
among them per request. Their $\mathbf{B}$ matrices were zero, so every output
had to equal the base model's, and every output did, token for token, while the
server still performed all the adapter arithmetic. One stream decoded at 41 to
42~ms per token with no adapter, 45 to 51~ms with a rank-16 adapter and 49 to
51~ms with rank 64 (measured, two runs each; record
\texttt{2026-09-24-m1-part3}).

Then two users arrived at once with two \emph{different} rank-16 adapters, and
the server could not put them in the same step. One user's 128 tokens were
complete after 6.7 seconds; the other's started only then and finished after
13.1. With the \emph{same} adapter, the two shared every step and both finished
after 11.8 seconds. The rule is in the server's source: two slots join the same
batch only if their tasks are of the same type and their adapter settings are
equal, and in each iteration the first decoding slot's settings decide
(\texttt{can\_batch\_with} and \texttt{update\_slots} in
\texttt{tools/server/server-context.cpp} at commit \texttt{d006858}). The
server held three adapters and served one at a time. Holding a thousand
fine-tunes is easy; this chapter is about serving them together.

\Foundations{A low-rank update is two smaller maps}{
An adapter can represent a weight change as a product of two thin matrices rather than
a full matrix. With row-vector convention,
$\mathbf{x}(\mathbf{W}+\mathbf{A}\mathbf{B})=\mathbf{xW}+(\mathbf{xA})\mathbf{B}$. The
inner width is the rank budget, not the model width. Requests may share the base
weights while selecting different adapter parameters. Sharing the base therefore does
not imply their cache entries are interchangeable. Chapter~\ref{ch:forward-pass}
explains the linear map; Chapter~\ref{ch:prefix-caching} explains the identity a
reusable prefix must carry.
}

\NapkinSection{Napkin math: an adapter's bytes and FLOPs}

A LoRA adapter replaces a weight matrix $\mathbf{W}\in\mathbb{R}^{d_{\text{out}}\times d_{\text{in}}}$
with $\mathbf{W} + \frac{\alpha}{r}\mathbf{B}\mathbf{A}$, where
$\mathbf{A}\in\mathbb{R}^{r\times d_{\text{in}}}$ and
$\mathbf{B}\in\mathbb{R}^{d_{\text{out}}\times r}$ \cite{hu2022lora}. Summed over
the adapted matrices of $L$ layers it adds
\begin{equation}
P_{\text{LoRA}} = r \, L \sum_{\text{matrices}} (d_{\text{in}} + d_{\text{out}})
\label{eq:lora-params}
\end{equation}
parameters, and $2P_{\text{LoRA}}$ FLOPs per token. For Llama~3.2~3B with all
seven projections adapted, the sum over one layer is 54,272, so
$P_{\text{LoRA}} = 1{,}519{,}616\,r$: 24.3 million at rank 16, 0.76\% of the
model's FLOPs per token (derived).

Bytes are a different story at decode. Each distinct adapter in a step must
be read once, like the base weights, so a batch that mixes $k$ adapters reads
the base model plus $k$ adapters. With the 2.01~GB quantized base of this book's
M1 measurements, forty-one different rank-16 adapters in one batch would read
as many adapter bytes as base-model bytes and roughly double the step time
(derived). Multi-adapter batching is cheap in arithmetic and paid for in
bandwidth, one adapter at a time.

The opening's single-stream cost fits neither estimate. The arithmetic was
under 1\% and the bytes 2.4\%, yet rank 16 added 2 to 10~ms per token. Each
adapted matrix adds four small operations to llama.cpp's graph --- two thin
matrix products, a scale and an addition (\texttt{build\_lora\_mm} in
\texttt{src/llama-graph.cpp} at commit \texttt{d006858}) --- 784 per step for
this model, and on this machine dispatching a small operation costs
microseconds (\Chref{kernel-languages}). That is a hypothesis the lab can
test; the measured cost is the number that matters.

\section{Low-rank adapters}

LoRA freezes the base model and trains only the low-rank factors. Its authors
reported ten thousand times fewer trainable parameters and a third of the GPU
memory of full fine-tuning for GPT-3 175B, with quality on par or better
\cite{hu2022lora}. For serving, the decisive property is that an adapter can be
merged into the base weights, $\mathbf{W}' = \mathbf{W} + \frac{\alpha}{r}\mathbf{B}\mathbf{A}$,
so that one fine-tune served alone costs nothing extra. Merging is exactly what
a shared server cannot do: a batch that mixes users of different adapters must
keep $\mathbf{W}$ shared and apply each adapter separately.

\section{Batching requests with different adapters}

With the base weights shared, a layer's output for a batch whose token $i$
uses adapter $a_i$ is
\begin{equation}
\mathbf{y}_i = \mathbf{W}\mathbf{x}_i + \tfrac{\alpha}{r}\,\mathbf{B}_{a_i}\left(\mathbf{A}_{a_i}\mathbf{x}_i\right).
\label{eq:multi-lora}
\end{equation}
The first term is one batched product for everyone, the lever of
\Chref{continuous-batching}. The second is a different small product for each
group of tokens that share an adapter. A server that cannot compute the second
term for mixed adapters in one step must split the batch by adapter, as the
opening measured; one that can keeps the whole batch together and pays only for
the small products.

\section{Gather-and-multiply kernels}

Punica's answer is a kernel that performs equation~\eqref{eq:multi-lora}'s
second term for a whole batch at once: tokens are grouped into segments by
adapter, and each segment gathers its adapter's factors and multiplies them,
all in one launch \cite{chen2024punica}. Its authors called it segmented
gather matrix--vector multiplication and reported twelve times the throughput
of serving systems of the time for multi-LoRA workloads, at a cost of about
2~ms of latency per token (their measurement). S-LoRA extended the kernels to
batches mixing adapters of different ranks \cite{sheng2024slora}.

\section{Adapter paging and caching}

A thousand adapters need not all be in accelerator memory. S-LoRA stores them
in host memory and fetches the ones the running requests use, and it manages
adapter weights and KV cache in one paged pool so that adapters of different
ranks and caches of different lengths share memory without fragmenting it;
its authors reported up to four times the throughput of the systems they
compared against and several orders of magnitude more adapters served
\cite{sheng2024slora} (their measurement). vLLM exposes the same tiers as
configuration: a number of adapters resident on the GPU, a maximum rank that
sizes their memory, and a host-memory cache, with endpoints to load and unload
adapters while the server runs \cite{vllm2026lora}.

\section{Routing between models}

Many models behind one endpoint raise a second question: which model should
answer? A router can send easy requests to a small model and hard ones to a
large one. RouteLLM trained such routers on human preference data and reported
cutting cost by more than half in some cases without reducing response quality
\cite{ong2024routellm} (their measurement). Routing interacts with every
cache in this part of the book: a request routed to a replica that does not
hold its adapter pays a load, and one routed away from its prefix pays a
prefill (\Chref{prefix-caching}).

\section{Multiplexing models and cold starts}

When different base models share hardware, the cost of switching is loading
weights. On the M1, Ollama answered an eight-token request in 0.28 seconds
when the model was already loaded. Alternating between two 3-billion-parameter
models, each request first had to load its model: 1.6 to 2.3 seconds of
loading, and 2.2 to 3.0 seconds per request in all (measured; the files were
likely in the operating system's cache, so this is close to the fastest a
switch can be here). A switch cost about eight answers. ServerlessLLM attacked
this cost with checkpoint formats and loaders
built for the storage hierarchy of GPU servers, live migration and
locality-aware placement, and reported 10 to 200 times lower latency than
serverless baselines \cite{fu2024serverlessllm} (their measurement). Adapters
are the extreme case of the same idea: make the part that differs between
models small enough that switching costs almost nothing.

Sharing base weights does not erase the identity of a fine-tuned model. The selected adapter controls which update is added, and that identity must follow the request through scheduling and cache lookup. An adapter miss is therefore not permission to serve the base result silently.

\EngineFigure{sys-adapter-components}{Component layout of one adapter-aware linear operation. Routing chooses the adapter identity before evaluating its low-rank update. The shared base contribution and the selected update are added for the same token rows. This is not the entire decoder.}{fig:sys-adapter-components}

\LedgerSection

Serving many fine-tunes from one base model with batched low-rank adapters:

\begin{ConstraintLedger}
Memory capacity & buys & One base model for all variants: 48.6~MB per rank-16 adapter of Llama~3.2~3B instead of a 6.4~GB copy. \\
Compute & spends & $2P_{\text{LoRA}}$ FLOPs per token, 0.76\% at rank 16, plus a few small operations per adapted matrix. \\
Memory bandwidth & spends & Every distinct adapter in a step is read in full: 48.6~MB each at rank 16. \\
Latency & risks & Without multi-adapter kernels, users of different adapters take turns: 13.1~s instead of 11.8~s for the second of two users on the M1. \\
\end{ConstraintLedger}

\LabSection{a batched multi-adapter GEMM}

\LabIntent{In \texttt{code/labs/ch21-multi-lora/}, a Rust implementation of
equation~\eqref{eq:multi-lora} three ways: merged weights per adapter, a loop
that splits the batch by adapter, and a segmented gather-and-multiply that
serves the whole batch in one pass. A float64 oracle checks all three. The
reader measures throughput as the number of distinct adapters in a 64-token
batch grows from 1 to 64, and finds where the gathered adapter bytes overtake
the base weights.}

\HappenedSection{multi-adapter serving in production}

Multi-adapter serving went from two research systems to a standard engine
feature. Punica and S-LoRA, both published in late 2023, showed that batching
across adapters needs its own kernels and its own memory management
\cite{chen2024punica,sheng2024slora}. vLLM now serves adapter requests in
the same batches as base-model requests, up to a configured number of
distinct adapters, and loads adapters at run time \cite{vllm2026lora}.
llama.cpp's server, in the build measured here, applies adapters per request
but batches only requests with identical adapter settings; the opening
measured what that means when two users differ. Neither design is wrong for its
setting: a laptop serving one user rarely mixes adapters, and a hosted
fine-tuning service does little else.

\FrontierSection{2026-09-24}

Batched multi-adapter serving is standard in data-center engines
\cite{vllm2026lora}, with adapters paged between host and GPU memory
\cite{sheng2024slora}. The frontier is at the fleet level: routing among many
models and adapters by cost and quality \cite{ong2024routellm}, and loading
models fast enough that multiplexing hardware between them is practical
\cite{fu2024serverlessllm}.

\ExercisesSection

\begin{enumerate}
\item \emph{Derivation.} Using equation~\eqref{eq:lora-params}, compute the
size of a rank-8 adapter on only the query and value projections of
Qwen3-8B, and the number of distinct such adapters a decode step can read
before adapter bytes equal the BF16 base weights.
\item \emph{Prediction.} In the opening's two-user experiment, predict the
completion times if the second user had asked for 512 tokens instead of 128.
\item \emph{Implementation.} In the lab, add adapters of different ranks to one
batch and compare padding every adapter to the largest rank with grouping by
rank.
\item \emph{Falsification.} Find a batch composition for which merging each
adapter into its own copy of the weights and running the groups one after
another is faster than a gather-and-multiply kernel.
\end{enumerate}
